% Source: https://artofproblemsolving.com/community/c3h1878211_volume_1_a_question_about_a_simple_thing_in_rationalizing_denominators
% Stable identity: aops-topic-1878211-q01
(a) Multiplying a product by $\sqrt[3]2$ multiplies only \emph{one} of its factors by $\sqrt[3]2$:
\[
\sqrt[3]{12}\cdot\sqrt[3]2=\left(\sqrt[3]{2^2}\cdot\sqrt[3]3\right)\sqrt[3]2=\left(\sqrt[3]{2^2}\cdot\sqrt[3]2\right)\sqrt[3]3=\sqrt[3]{8}\cdot\sqrt[3]3=2\sqrt[3]3 .
\]
Writing $2\sqrt[3]6$ would multiply by $\sqrt[3]2$ twice (once to turn $\sqrt[3]{2^2}$ into $2$ and again to turn $\sqrt[3]3$ into $\sqrt[3]6$). Directly: $\sqrt[3]{12}\cdot\sqrt[3]2=\sqrt[3]{24}=\sqrt[3]{8\cdot3}=2\sqrt[3]3$.

(b) The result $\dfrac{\sqrt[3]2}{2\sqrt[3]3}$ still has $\sqrt[3]3$ in the denominator. To make $12=2^2\cdot3$ into a perfect cube we need one more $2$ and two more $3$s, so multiply by $\sqrt[3]{2\cdot3^2}=\sqrt[3]{18}$:
\[
\dfrac{1}{\sqrt[3]{12}}=\dfrac{\sqrt[3]{18}}{\sqrt[3]{216}}=\boxed{\dfrac{\sqrt[3]{18}}{6}}.
\]
